% ============================================================
\section{符号说明}
\begin{table}[H]
  \centering
  \label{tab:notation}
  \footnotesize
  \renewcommand{\arraystretch}{1.05}
  \begin{tabular}{cll}
    \toprule
    符号 & 说明 & 单位 \\
    \midrule
    $t$ & 10分钟时段编号，$t=1,\ldots,144$ & -- \\
    $d$ & 日期编号（滚动预测/调度的当前天） & -- \\
    $p_t$ & 时段 $t$ 的购电电价 & 元/kWh \\
    $L_t$ & 时段 $t$ 的负载电量 & kWh \\
    $V_t$ & 时段 $t$ 的光伏预测电量 & kWh \\
    $g_t,\,g_t^0$ & 时段 $t$ 的（原）计划购电量 & kWh \\
    $c_t,\,d_t$ & 储能设备外部端口的充电量、放电量 & kWh \\
    $w_t$ & 未利用（弃用）的光伏电量 & kWh \\
    $E_t,\,E_0,\,E_{\mathrm{end}}$ & 时段 $t$末、初始（0:00）、当天24:00的储能电量 & kWh \\
    $\eta_c,\,\eta_d$ & 储能充电效率、放电效率 & 无量纲 \\
    $z_t$ & 储能充电模式二元变量（1为充电、0为放电） & 0/1 \\
    $M_c,\,M_d$ & 单时段最大充电电量、最大放电电量 & kWh \\
    $F,\,F_{\mathrm{LP}}^\ast,\,F_{\mathrm{MILP}}^\ast$ & 全天购电费用（目标函数）、LP松弛与原MILP最优值 & 元 \\
    $u_t,\,v_t$ & 日内调整中购电量的追加量、减购量 & kWh \\
    $h_t,\,h_t^{(j)}$ & 实际（或按情景 $j$ 计）发生的紧急购电量 & kWh \\
    $a_t,\,a_t^{(j)}$ & 净负荷电量、按历史误差情景 $j$ 构造的净负荷情景 & kWh \\
    $C,\,C_{\mathrm{purchase}},\,C_{\mathrm{opt}}$ & 全年总费用、购电类费用合计、带惩罚项的优化目标 & 元 \\
    $W$ & 净负荷预测模型的滚动训练窗口长度 & 天 \\
    $K_L,\,K_P$ & 负荷、光伏PCA保留的主成分个数 & -- \\
    $\widehat{\bm L}_d,\,\widehat{\bm P}_d$ & 第 $d$ 天负荷、光伏曲线的中心预测 & kWh \\
    $q,\,q^\ast$ & 计划购电所依据的净负荷风险分位数、开发期扫描确定的最优值 & -- \\
    $\tau$ & 光伏预报发布时刻（00:00/06:00/12:00/18:00） & 时:分 \\
    $x$ & 期末可用储能量（以容量下限1200\,kWh为基准） & kWh \\
    \bottomrule
  \end{tabular}
\end{table}
